\renewcommand{\bibname}{Quellen und Bandzuordnung}
\begin{thebibliography}{99}
\addcontentsline{toc}{chapter}{Quellen und Bandzuordnung}
\footnotesize
\setlength{\itemsep}{5pt}
\setlength{\parskip}{0pt}
\interlinepenalty=10000

\bibitem{project}
Martin Kunze: \emph{Die Grundlagen der Mathematik}, Band~10:
\emph{Natürliche Zahlen}.
Grundlage für Konstruktion, Peano-Axiome, Nachfolger, Induktion,
Dedekindschen Rekursionssatz, Addition, natürliche Ordnung,
Multiplikation, endliche Summen und Produkte, Potenzen, Fakultäten
und Stellenwertnotation. Die Leitfäden der Bände~3, 4, 5, 28
und~37 dienen als Formatvorbilder.
Die vorliegende Ergänzung enthält zusätzlich ausgewählte
Standardfolgen, insbesondere Binomialformeln, Division mit Rest,
elementare Teilbarkeit, Primexponenten, ausgewählte klassische
arithmetische Hauptsätze, Teileranzahl, Teilersumme und Eulers
$\varphi$-Funktion sowie die Eindeutigkeit
der Stellenwertdarstellung. Damit wird nicht behauptet, daß
jeder dieser Sätze bereits im Hauptband einzeln bewiesen ist.

\bibitem{dedekind}
Richard Dedekind: \emph{Was sind und was sollen die Zahlen?}, 1888.
Englische Übersetzung von Wooster Woodruff Beman in
\emph{Essays on the Theory of Numbers}, Open Court, 1901;
zweite Abhandlung \emph{The Nature and Meaning of Numbers}.
Nr.~71 und~80: einfach unendliches System und Induktion;
Nr.~126 und~132: Rekursion und Eindeutigkeit des Zahlensystems;
Abschnitte~XI--XIII: Addition, Multiplikation, Potenzen.
Dedekinds Anfangselement heißt $1$; hier wird mit $0$ begonnen.
\url{https://www.gutenberg.org/files/21016/21016-pdf.pdf}.

\bibitem{natchoose}
The mathlib Community: \emph{Mathlib.Data.Nat.Choose.Basic},
offizielle Dokumentation, abgerufen am 24.~Oktober 2026.
Vergleich für natürliche Binomialkoeffizienten:
\path{Nat.choose}, \path{Nat.choose_succ_succ},
\path{Nat.choose_mul_factorial_mul_factorial}, \path{Nat.choose_symm}.
\url{https://leanprover-community.github.io/mathlib4_docs/Mathlib/Data/Nat/Choose/Basic.html}.

\bibitem{natgcd}
The mathlib Community: \emph{Mathlib.Data.Nat.GCD.Basic} und
die dort verlinkten Definitionen von \path{Nat.gcd},
\path{Nat.lcm} und \path{Nat.Coprime}, offizielle Dokumentation,
abgerufen am 24.~Oktober 2026.
Abgleich der Teilbarkeitscharakterisierungen und teilerfremden
Kürzungsregeln, insbesondere \path{Nat.gcd_greatest},
\path{Nat.Coprime.dvd_mul_right} und \path{Nat.Coprime.lcm_eq_mul}.
\path{Nat.pow_sub_one_gcd_pow_sub_one} liefert außerdem den ggT
zweier Potenzdifferenzen zur gleichen natürlichen Basis.
\url{https://leanprover-community.github.io/mathlib4_docs/Mathlib/Data/Nat/GCD/Basic.html}.

\bibitem{natprime}
The mathlib Community: \emph{Mathlib.Data.Nat.Prime.Defs} und
\emph{Mathlib.Data.Nat.Prime.Basic}, offizielle Dokumentation,
abgerufen am 24.~Oktober 2026.
Vergleich für Primzahlbegriff, Primteiler und Produktkriterien:
\path{Nat.prime_def}, \path{Nat.prime_def_le_sqrt},
\path{Nat.prime_iff_not_exists_mul_eq},
\path{Nat.ne_one_iff_exists_prime_dvd}.
\url{https://leanprover-community.github.io/mathlib4_docs/Mathlib/Data/Nat/Prime/Defs.html};
\url{https://leanprover-community.github.io/mathlib4_docs/Mathlib/Data/Nat/Prime/Basic.html}.

\bibitem{natfactors}
The mathlib Community: \emph{Mathlib.Data.Nat.Factors},
offizielle Dokumentation, abgerufen am 24.~Oktober 2026.
Vergleich für die Primfaktorzerlegung:
\path{Nat.prod_primeFactorsList}, \path{Nat.primeFactorsList_sorted}
und \path{Nat.primeFactorsList_unique}.
\url{https://leanprover-community.github.io/mathlib4_docs/Mathlib/Data/Nat/Factors.html}.

\bibitem{natinfinite}
The mathlib Community: \emph{Mathlib.Data.Nat.Prime.Infinite},
offizielle Dokumentation, abgerufen am 24.~Oktober 2026.
\path{Nat.exists_infinite_primes}: Zu jeder natürlichen Schranke
gibt es eine mindestens so große Primzahl.
\url{https://leanprover-community.github.io/mathlib4_docs/Mathlib/Data/Nat/Prime/Infinite.html}.

\bibitem{natdigits}
The mathlib Community: \emph{Mathlib.Data.Nat.Digits.Defs},
offizielle Dokumentation, abgerufen am 24.~Oktober 2026.
Vergleich für Auswertung, Stellenzahl und Rückgewinnung von Ziffern:
\path{Nat.ofDigits_digits}, \path{Nat.ofDigits_inj_of_len_eq},
\path{Nat.ofDigits_lt_base_pow_length},
\path{Nat.self_div_pow_eq_ofDigits_drop} und
\path{Nat.self_mod_pow_eq_ofDigits_take}.
Dort wird die Null auch durch eine leere Ziffernliste behandelt;
hier gilt die im Hauptband verwendete Normalform mit einer Nullziffer.
\url{https://leanprover-community.github.io/mathlib4_docs/Mathlib/Data/Nat/Digits/Defs.html}.

\bibitem{natvandermonde}
The mathlib Community: \emph{Mathlib.Data.Nat.Choose.Vandermonde},
offizielle Dokumentation, abgerufen am 24.~Oktober 2026.
\path{Nat.add_choose_eq} und \path{Nat.sum_range_choose_sq}:
Vandermonde-Identität und Quadratsumme der Binomialkoeffizienten.
\url{https://leanprover-community.github.io/mathlib4_docs/Mathlib/Data/Nat/Choose/Vandermonde.html}.

\bibitem{natgeomsum}
The mathlib Community: \emph{Mathlib.Algebra.Ring.GeomSum},
offizielle Dokumentation, abgerufen am 24.~Oktober 2026.
\texttt{geom\_sum\textsubscript{2}\_mul\_of\_ge}: Faktorisierung einer Potenzdifferenz.
Verwendet wird ausschließlich die Fassung für natürliche Zahlen
mit größerer oder gleicher erster Basis.
\url{https://leanprover-community.github.io/mathlib4_docs/Mathlib/Algebra/Ring/GeomSum.html}.

\bibitem{natroot}
The mathlib Community:
\emph{Mathlib.Analysis.SpecialFunctions.Pow.NthRootLemmas},
offizielle Dokumentation, abgerufen am 24.~Oktober 2026.
\path{Nat.pow_nthRoot_le}, \path{Nat.lt_pow_nthRoot_add_one},
\path{Nat.le_nthRoot_iff} und \path{Nat.nthRoot_pow}:
Wurzelschranken und Potenzerkennung für beliebige positive
natürliche Exponenten. Gemeint ist der natürliche Wurzelwert,
charakterisiert durch $q^r\le n<(q+1)^r$.
\url{https://leanprover-community.github.io/mathlib4_docs/Mathlib/Analysis/SpecialFunctions/Pow/NthRootLemmas.html}.

\bibitem{natdivmod}
The Lean Developers: \emph{Init.Data.Nat.Div.Basic} und
\emph{Init.Data.Nat.Lemmas}, offizielle Dokumentation,
abgerufen am 24.~Oktober 2026.
Vergleichsstellen: \path{Nat.add_mul_div_left},
\path{Nat.div_div_eq_div_mul}, \path{Nat.mod_mod},
\path{Nat.add_mod} und \path{Nat.mul_mod}.
\url{https://leanprover-community.github.io/mathlib4_docs/Init/Data/Nat/Div/Basic.html};
\url{https://leanprover-community.github.io/mathlib4_docs/Init/Data/Nat/Lemmas.html}.

\bibitem{natvaluation}
The mathlib Community: \emph{Mathlib.Data.Nat.Factorization.Defs}
und \emph{Mathlib.Data.Nat.Factorization.Basic}, offizielle
Dokumentation, abgerufen am 24.~Oktober 2026.
Vergleich für Primexponenten:
\path{Nat.factorization_mul}, \path{Nat.factorization_pow},
\path{Nat.factorization_le_iff_dvd}, \path{Nat.eq_of_factorization_eq},
\path{Nat.factorization_gcd} und \path{Nat.factorization_lcm}.
Der Leitfaden verwendet Primexponenten nur für positive Zahlen.
\url{https://leanprover-community.github.io/mathlib4_docs/Mathlib/Data/Nat/Factorization/Defs.html};
\url{https://leanprover-community.github.io/mathlib4_docs/Mathlib/Data/Nat/Factorization/Basic.html}.

\bibitem{natlegendre}
The mathlib Community: \emph{Mathlib.NumberTheory.Padics.PadicVal.Basic},
offizielle Dokumentation, abgerufen am 24.~Oktober 2026.
\path{padicValNat_factorial}: Legendre-Formel als endliche Summe
natürlicher Quotienten. Hier wird die Summe bis zum Exponenten $n$
genommen; alle höheren Beiträge sind $0$.
\path{padicValNat_choose},
\path{sub_one_mul_padicValNat_choose_eq_sub_sum_digits} und
\path{sub_one_mul_padicValNat_factorial} dienen zusätzlich dem
Abgleich von Kummers Übertragssatz und den Ziffernsummenformeln.
Alle Differenzformeln sind hier als Gleichungen natürlicher
Summen umgeschrieben.
\url{https://leanprover-community.github.io/mathlib4_docs/Mathlib/NumberTheory/Padics/PadicVal/Basic.html}.

\bibitem{natfermat}
The mathlib Community: \emph{Mathlib.FieldTheory.Finite.Basic},
offizielle Dokumentation, abgerufen am 24.~Oktober 2026.
\path{Nat.ModEq.pow_card_sub_one_eq_one} und
\path{Nat.ModEq.pow_totient}: Fermats kleiner Satz und Eulers Satz
für natürliche Zahlen. Beide werden hier als Gleichheit
natürlicher Restwerte formuliert.
\url{https://leanprover-community.github.io/mathlib4_docs/Mathlib/FieldTheory/Finite/Basic.html}.

\bibitem{natwilson}
The mathlib Community: \emph{Mathlib.NumberTheory.Wilson},
offizielle Dokumentation, abgerufen am 24.~Oktober 2026.
\path{Nat.prime_iff_fac_equiv_neg_one}: Wilsons Primzahlkriterium;
hier als natürliche Teilbarkeitsbedingung für $(n-1)!+1$ formuliert.
\url{https://leanprover-community.github.io/mathlib4_docs/Mathlib/NumberTheory/Wilson.html}.

\bibitem{natchinese}
The mathlib Community: \emph{Mathlib.Data.Nat.ModEq},
offizielle Dokumentation, abgerufen am 24.~Oktober 2026.
\path{Nat.chineseRemainder}, \path{Nat.chineseRemainder_lt_mul}
und \path{Nat.chineseRemainder_modEq_unique}:
Existenz, Schranke und Eindeutigkeit beim chinesischen Restsatz.
Die Fassung ohne Teilerfremdheit entspricht
\path{Nat.chineseRemainder'}, \path{Nat.chineseRemainder'_lt_lcm}
und \path{Nat.modEq_and_modEq_iff_modEq_lcm}.
Im Leitfaden werden nur natürliche Restwerte und ihre Lösungen benutzt.
\url{https://leanprover-community.github.io/mathlib4_docs/Mathlib/Data/Nat/ModEq.html}.

\bibitem{natbertrand}
The mathlib Community: \emph{Mathlib.NumberTheory.Bertrand},
offizielle Dokumentation, abgerufen am 24.~Oktober 2026.
\path{Nat.exists_prime_lt_and_le_two_mul}: Zu jedem $n>0$ gibt es
eine Primzahl $p$ mit $n<p\le2n$. Für $n\ge2$ ist die obere
Ungleichung strikt.
\url{https://leanprover-community.github.io/mathlib4_docs/Mathlib/NumberTheory/Bertrand.html}.

\bibitem{natsquares}
The mathlib Community: \emph{Mathlib.NumberTheory.SumTwoSquares}
und \emph{Mathlib.NumberTheory.SumFourSquares}, offizielle
Dokumentation, abgerufen am 24.~Oktober 2026.
\path{Nat.eq_sq_add_sq_iff} und \path{Nat.Prime.sq_add_sq}:
Kriterium für zwei Quadrate und Primzahl-Spezialfall;
\path{Nat.sum_four_squares}: Lagranges Vier-Quadrate-Satz.
Es werden ausschließlich Darstellungen durch natürliche Quadrate
aufgenommen.
\url{https://leanprover-community.github.io/mathlib4_docs/Mathlib/NumberTheory/SumTwoSquares.html};
\url{https://leanprover-community.github.io/mathlib4_docs/Mathlib/NumberTheory/SumFourSquares.html}.


\bibitem{natlucas}
The mathlib Community: \emph{Mathlib.Data.Nat.Choose.Lucas},
offizielle Dokumentation, abgerufen am 24.~Oktober 2026.
\path{Choose.lucas_theorem_nat}: Berechnung des Restes eines
Binomialkoeffizienten aus den Ziffern zur Primzahlbasis.
\url{https://leanprover-community.github.io/mathlib4_docs/Mathlib/Data/Nat/Choose/Lucas.html}.

\bibitem{natlinearrest}
The Isabelle Developers: \emph{HOL-Number\_Theory.Pocklington},
offizielle Dokumentation, abgerufen am 24.~Oktober 2026.
\path{cong_solve} und \path{cong_solve_unique}: Lösbarkeit und
Eindeutigkeit der reduzierten natürlichen Restgleichung.
Zusammen mit \path{Nat.ModEq.cancel_left_div_gcd} aus
\cite{natchinese} dient dies dem Abgleich der hier formulierten
allgemeinen ggT-Bedingung und Lösungsschrittweite.
\path{ord_divides}, \path{order_divides_totient},
\path{order_divides_expdiff} und \path{ord_power} behandeln die
hier ebenfalls verwendeten natürlichen Rückkehr-Exponenten
von Potenzresten.
\url{https://www.isabelle.in.tum.de/library/HOL/HOL-Number_Theory/Pocklington.html}.

\bibitem{natfrobenius}
The mathlib Community: \emph{Mathlib.NumberTheory.FrobeniusNumber},
offizielle Dokumentation, abgerufen am 24.~Oktober 2026.
\path{frobeniusNumber_pair} und \path{frobeniusNumber_iff}:
größte nicht als $ax+by$ mit $x,y\in\mathbb N$ darstellbare Zahl
für zwei teilerfremde natürliche Koeffizienten $a,b\ge2$.
\url{https://leanprover-community.github.io/mathlib4_docs/Mathlib/NumberTheory/FrobeniusNumber.html}.

\bibitem{natpythagoras}
The mathlib Community: \emph{Mathlib.NumberTheory.PythagoreanTriples},
offizielle Dokumentation, abgerufen am 24.~Oktober 2026.
\path{PythagoreanTriple.coprime_classification} und
\path{PythagoreanTriple.coprime_classification'} sowie die dortige
Normalisierung. Der Leitfaden übernimmt ausschließlich die
Spezialisierung auf positive natürliche Lösungen und Parameter.
\url{https://leanprover-community.github.io/mathlib4_docs/Mathlib/NumberTheory/PythagoreanTriples.html}.

\bibitem{natthreesquares}
Anna Danilkin und Loïc Chevalier: \emph{Three Squares Theorem},
Archive of Formal Proofs, 3.~Mai 2023; abgerufen am 24.~Oktober 2026.
Vergleich für das Drei-Quadrate-Kriterium in seiner Formulierung
mit natürlichen Zahlen und natürlichen Quadratwurzeln der Summanden.
\url{https://isa-afp.org/entries/Three_Squares.html}.

\bibitem{natarithfunctions}
The mathlib Community: \emph{Mathlib.NumberTheory.ArithmeticFunction.Misc},
offizielle Dokumentation, abgerufen am 24.~Oktober 2026.
\path{Nat.card_divisors}, \path{Nat.sum_divisors},
\path{Nat.Coprime.card_divisors_mul} und
\path{Nat.Coprime.sum_divisors_mul}: Teileranzahl, Teilersumme
und Produktregeln bei teilerfremden Argumenten.
\url{https://leanprover-community.github.io/mathlib4_docs/Mathlib/NumberTheory/ArithmeticFunction/Misc.html}.

\bibitem{nattotient}
The mathlib Community: \emph{Mathlib.Data.Nat.Totient},
offizielle Dokumentation, abgerufen am 24.~Oktober 2026.
\path{Nat.totient_prime_pow}, \path{Nat.totient_eq_prod_factorization},
\path{Nat.totient_mul} und \path{Nat.sum_totient}:
Primzahlpotenzformel, Produktformeln und Teilersumme der Euler-Funktion.
Im Leitfaden ist ihr Argument stets positiv.
\url{https://leanprover-community.github.io/mathlib4_docs/Mathlib/Data/Nat/Totient.html}.


\bibitem{natvalsum}
The mathlib Community: \emph{Mathlib.RingTheory.Multiplicity},
offizielle Dokumentation, abgerufen am 24.~Oktober 2026.
\path{min_le_emultiplicity_add} und \path{multiplicity_add_eq_min}:
Vergleich der Primexponenten bei Addition. Verwendet wird nur
die Spezialisierung auf positive natürliche Summanden und Primzahlen.
\url{https://leanprover-community.github.io/mathlib4_docs/Mathlib/RingTheory/Multiplicity.html}.

\bibitem{natbezout}
The Isabelle Developers: \emph{HOL.GCD}, offizielle Dokumentation,
abgerufen am 24.~Oktober 2026.
\path{bezout_nat}: Bézouts Aussage mit natürlichen Koeffizienten
in der Form $ax=by+\operatorname{ggT}(a,b)$ für $a>0$.
\url{https://www.isabelle.in.tum.de/library/HOL/HOL/GCD.html}.

\bibitem{natdirichlet}
The mathlib Community: \emph{Mathlib.NumberTheory.LSeries.PrimesInAP},
offizielle Dokumentation, abgerufen am 24.~Oktober 2026.
\path{Nat.forall_exists_prime_gt_and_modEq}: Dirichlets Satz
mit beliebiger natürlicher unterer Schranke und teilerfremder Restvorgabe.
Der Leitfaden verwendet nur die Aussage über natürliche Primzahlen.
\url{https://leanprover-community.github.io/mathlib4_docs/Mathlib/NumberTheory/LSeries/PrimesInAP.html}.

\bibitem{natperfect}
The mathlib Community: \emph{Archive.Wiedijk100Theorems.PerfectNumbers},
offizielle Dokumentation, abgerufen am 24.~Oktober 2026.
\path{Theorems100.Nat.even_and_perfect_iff}: Euklid--Euler-Kriterium
für gerade vollkommene natürliche Zahlen. Der Zusammenhang mit der
Teilersumme ist in \emph{Mathlib.NumberTheory.Divisors} als
\path{Nat.perfect_iff_sum_divisors_eq_two_mul} angegeben.
\url{https://leanprover-community.github.io/mathlib4_docs/Archive/Wiedijk100Theorems/PerfectNumbers.html};
\url{https://leanprover-community.github.io/mathlib4_docs/Mathlib/NumberTheory/Divisors.html}.

\bibitem{natprimroots}
Manuel Eberl / The Isabelle Developers:
\emph{HOL-Number\_Theory.Residue\_Primitive\_Roots},
offizielle Dokumentation, abgerufen am 24.~Oktober 2026.
\path{residue_primroot_is_generator},
\path{residue_primroot_iff_in_cyclic_moduli} und
\path{card_residue_primroots}: Darstellung teilerfremder Restwerte,
Existenz und Anzahl primitiver Wurzeln.
Die Aussagen sind hier ausschließlich in natürlicher Restwertnotation
formuliert.
\url{https://www.isabelle.in.tum.de/library/HOL/HOL-Number_Theory/Residue_Primitive_Roots.html}.


\bibitem{natlte}
The mathlib Community: \emph{Mathlib.NumberTheory.Multiplicity},
offizielle Dokumentation, abgerufen am 24.~Oktober 2026.
\path{padicValNat.pow_sub_pow} und \path{padicValNat.pow_two_sub_pow}:
Primexponenten von Potenzdifferenzen für ungerade Primzahlen und für $2$.
Im zweiten Fall ist die Formel hier als Gleichheit natürlicher Summen
formuliert.
\url{https://leanprover-community.github.io/mathlib4_docs/Mathlib/NumberTheory/Multiplicity.html}.

\bibitem{natfermatlast}
Andrew Wiles: \emph{Modular elliptic curves and Fermat's Last Theorem},
Annals of Mathematics, Band~141, Nr.~3 (1995), S.~443--551;
Satz~0.5 auf S.~448. Hier wird nur die Aussage über natürliche
Lösungen von $x^r+y^r=z^r$ übernommen.
\url{https://annals.math.princeton.edu/1995/141-3/p01};
\href{https://wstein.org/edu/Fall2003/252/references/wiles-flt/Wiles-Modular_Elliptic_Curves_and_Fermats_Last_Theorem.pdf}{Digitalisat des Artikels}.

\bibitem{natcarmichael}
The mathlib Community: \emph{Mathlib.NumberTheory.ArithmeticFunction.Carmichael},
offizielle Dokumentation, abgerufen am 24.~Oktober 2026.
\path{ArithmeticFunction.carmichael_dvd_totient},
\path{ArithmeticFunction.carmichael_factorization},
\path{ArithmeticFunction.carmichael_two_pow_of_le_two},
\path{ArithmeticFunction.carmichael_two_pow_of_ne_two} und
\path{ArithmeticFunction.carmichael_pow_of_prime_ne_two}:
Teilbarkeits- und Berechnungsformeln für den universellen Exponenten.
Die Charakterisierung und Realisierung als Rückkehr-Exponent
stehen außerdem als \path{Carmichael_divides} und
\path{Carmichael_root_exists} in \cite{natprimroots}.
\url{https://leanprover-community.github.io/mathlib4_docs/Mathlib/NumberTheory/ArithmeticFunction/Carmichael.html}.

\bibitem{natkorselt}
The mathlib Community: \emph{Mathlib.NumberTheory.CarmichaelNumber},
offizielle Dokumentation, abgerufen am 24.~Oktober 2026.
\path{Nat.isCarmichael_iff_korselt}: Charakterisierung der
Carmichael-Zahlen durch Primteiler und deren Exponenten.
\url{https://leanprover-community.github.io/mathlib4_docs/Mathlib/NumberTheory/CarmichaelNumber.html}.

\bibitem{natquadraticbasic}
The mathlib Community: \emph{Mathlib.NumberTheory.LegendreSymbol.Basic},
offizielle Dokumentation, abgerufen am 24.~Oktober 2026.
\path{ZMod.euler_criterion}, \path{ZMod.pow_div_two_eq_neg_one_or_one},
\path{legendreSym.card_sqrts} und \path{ZMod.exists_sq_eq_neg_one_iff}:
Euler-Kriterium, Anzahl der Quadratwurzeln und erster Ergänzungssatz.
Die Aussagen werden hier mit natürlichen Restwerten, insbesondere
$p-1$, statt mit zusätzlichen Zahlbereichen formuliert.
\url{https://leanprover-community.github.io/mathlib4_docs/Mathlib/NumberTheory/LegendreSymbol/Basic.html}.

\bibitem{natquadraticreciprocity}
The mathlib Community:
\emph{Mathlib.NumberTheory.LegendreSymbol.QuadraticReciprocity},
offizielle Dokumentation, abgerufen am 24.~Oktober 2026.
\path{ZMod.exists_sq_eq_two_iff},
\path{ZMod.exists_sq_eq_prime_iff_of_mod_four_eq_one} und
\path{ZMod.exists_sq_eq_prime_iff_of_mod_four_eq_three}:
zweiter Ergänzungssatz und quadratisches Reziprozitätsgesetz
als Aussagen über die Lösbarkeit von Quadratgleichungen.
\url{https://leanprover-community.github.io/mathlib4_docs/Mathlib/NumberTheory/LegendreSymbol/QuadraticReciprocity.html}.

\end{thebibliography}

\noindent\textbf{Umfang des Quellenabgleichs.}
Die Quellen dienen der Kontrolle der Begriffe und ausgewählter
Standardresultate. Zusammengefaßte Aussagen sind für diesen Leitfaden
in einheitlicher Notation formuliert; die Quellen sind keine
Behauptung wortgleicher Formulierungen oder historischer Erstentdeckungen.
Die Darstellung ist eine Auswahl zentraler Resultate, kein vollständiger
Katalog der Zahlentheorie. Auch die Verwendung formalisierter
Vergleichsquellen bedeutet keine maschinelle Formalverifikation
dieses Manuskripts.

